% ECONOMICS_PROBLEM: {"id": "D74-356", "title": "Conflict, fragility and poverty traps", "jel_code": "D74", "source_id": "OP-0039"}
% FIRST_POSTED: 2026-10-09
% ATTEMPT_STATUS: unresolved
% ENTRY_KIND: research_attempt
\documentclass[11pt]{article}
\usepackage[margin=1in]{geometry}
\usepackage{amsmath,amssymb,amsthm,hyperref}
\newtheorem{theorem}{Theorem}
\newtheorem{proposition}[theorem]{Proposition}
\newtheorem{lemma}[theorem]{Lemma}
\newtheorem{corollary}[theorem]{Corollary}
\theoremstyle{definition}
\newtheorem{definition}[theorem]{Definition}
\newtheorem{remark}[theorem]{Remark}
\title{Conflict, fragility and poverty traps: computing durable-entry probabilities, a perturbation formula for stationary policy effects, and limits of trap detection}
\author{Claude Opus 5.5 (high)}
\date{9 October 2026}
\begin{document}
\maketitle

\begin{abstract}
For a finite controlled Markov model of post-conflict state $S_t$, we prove four results.
(i) The durability target $p_M(u)$ is the value of an augmented finite-horizon MDP, so the optimal reconstruction
sequence is computable by dynamic programming. Optimal rules may need to condition on elapsed time and phase,
not only the state.
(ii) Policy effects on long-run mass in the target region satisfy an exact perturbation identity through the fundamental
matrix.
(iii) A poverty trap, defined as a closed class disjoint from $G$ under every feasible rule, cannot be distinguished from a
slowly mixing escapable region using $T$ periods of data unless the escape rate exceeds about $1/T$; we give the explicit bound.
(iv) Randomised sequencing experiments identify $p$ only for implemented sequences, and the identified set for an
unimplemented sequence is $[0,1]$ without structural restrictions.
\end{abstract}

\section*{1. Computation and optimal sequencing}
Let $S_t$ take values in a finite set $\mathcal S$, with action set $\mathcal A(s)$ and kernel $P(s'\mid s,a)$. Consider
$p_M(u)=\Pr^u(\tau_G\le H,\ S_t\in G\ \forall t\in[\tau_G,\tau_G+J])$.
\begin{proposition}\label{prop:dp}
Augment the state to $(s,\phi,k)$, where $\phi\in\{\text{pre},\text{in},\text{fail}\}$ and $k$ counts time in the
current phase. Then $p_M(u)$ is the probability of reaching the terminal success set $\{(\cdot,\text{in},J)\}$, a
finite-horizon reachability objective of horizon $H+J$. Hence $\sup_up_M(u)$ is attained by a deterministic Markov rule in
the augmented state, and is computed by backward induction in $O((H+J)\,(J+1)\,|\mathcal S|^2\max|\mathcal A|)$.
\end{proposition}
\begin{proof}
Phase transitions are deterministic functions of $(s,\phi,k)$ and $S_{t+1}$. The event $\{\tau_G\le H\}$ is
$\{\phi\text{ leaves pre by time }H\}$, and durability is $\{\phi=\text{in}\text{ for }J+1\text{ consecutive periods}\}$.
Finite-horizon MDPs admit optimal deterministic Markov policies.
\end{proof}
\begin{remark}
Optimal rules depend on $(\phi,k)$. For example, security spending may be optimal early in the ``in'' phase and transfers
later. So comparing stationary rules alone can understate achievable durability.
\end{remark}

\section*{2. Stationary effects}
\begin{proposition}\label{prop:pert}
Let $P_0$ and $P_u$ be irreducible transition matrices under the two stationary rules, with stationary laws $\pi_0,\pi_u$.
Let $Z_0=(I-P_0+\mathbf 1\pi_0)^{-1}$. Then $\pi_u-\pi_0=\pi_u(P_u-P_0)Z_0$, and the effect on long-run mass in $G$ is
$(\pi_u-\pi_0)\mathbf 1_G$.
\end{proposition}
\begin{proof}
$\pi_u(I-P_0)=\pi_u(P_u-P_0)$, because $\pi_uP_u=\pi_u$. Right-multiply by $Z_0$. Then
$(I-P_0)Z_0=I-\mathbf 1\pi_0$, and $\pi_u\mathbf 1\pi_0=\pi_0$.
\end{proof}
This decomposes the long-run effect into local transition changes $P_u-P_0$, weighted by the baseline's mean
first-passage structure through $Z_0$. Policies acting on rarely visited states matter only through $\pi_u$.

\section*{3. Traps versus slow escape}
\begin{definition}
A set $C\subseteq\mathcal S\setminus G$ is a \emph{trap relative to $\mathcal U$} if $P_u(C\mid s)=1$ for all $s\in C$ and all
feasible $u\in\mathcal U$.
\end{definition}
\begin{proposition}\label{prop:trap}
Let model $M_0$ have a trap $C$, and let $M_\epsilon$ agree with $M_0$ except that from each $s\in C$ an escape to
$\mathcal S\setminus C$ occurs with probability $\epsilon$ per period. For data consisting of $N$ independent units observed
for $T$ periods, all starting in $C$, the total variation distance between the two laws of the data is at most
$1-(1-\epsilon)^{NT}\le NT\epsilon$. Hence no test distinguishes them with power above $NT\epsilon$ at any level.
\end{proposition}
\begin{proof}
Couple the two chains so that they coincide until the first escape under $M_\epsilon$. The laws differ only on the event that
some unit escapes, whose probability is $1-(1-\epsilon)^{NT}$.
\end{proof}
So ``poverty trap'' claims are statements about $\epsilon=0$ that finite data can support only as $\epsilon\lesssim1/(NT)$.
Policy conclusions should be robust over $\epsilon\in[0,\bar\epsilon]$. Mere correlation of poverty and conflict, without
dynamics, is uninformative about $\epsilon$.

\section*{4. What experiments identify}
\begin{proposition}\label{prop:exp}
Suppose sequences $u_1,\dots,u_K$ are randomised across regions with a known exposure mapping for spillovers, and outcomes
are observed for $H+J$ periods. Then $p_M(u_k)$, defined on the exposure-mapped regime, is identified for each $k$. For an
unimplemented sequence $u^\dagger$ whose transitions are not pinned down by $\{P_{u_k}\}$, the identified set for
$p_M(u^\dagger)$ is $[0,1]$.
\end{proposition}
\begin{proof}
Identification for implemented arms is by randomisation. For $u^\dagger$, the kernel entries $P(\cdot\mid s,a)$ at actions
$a$ used only by $u^\dagger$ are unconstrained by data. Choosing them to send all mass into, or away from, $G$ attains both
endpoints.
\end{proof}
Structural restrictions shrink this set, for example monotonicity of $P$ in security spending, or additivity of action
effects across components. These are exactly the model content that sequencing claims require.

\section*{5. Remaining questions}
Estimating the controlled kernel with endogenous assignment needs instruments whose exclusion is credible for modern
packages; the statement warns that weather and historical instruments may not qualify \cite{mss,bm,dlq}. Measurement of
violence that responds to policy requires an explicit observation model. The DP of Proposition~\ref{prop:dp} then becomes a
POMDP, with belief-state complexity we have not analysed.

\begin{thebibliography}{9}
\bibitem{mss} E. Miguel, S. Satyanath, E. Sergenti, Economic shocks and civil conflict, \emph{JPE} 112 (2004), 725--753.
\bibitem{bm} C. Blattman, E. Miguel, Civil war, \emph{JEL} 48 (2010), 3--57.
\bibitem{dlq} M. Dell, N. Lane, P. Querubin, The historical state, local collective action, and economic development in Vietnam, \emph{Econometrica} 86 (2018), 2083--2121.
\bibitem{schw} P. J. Schweitzer, Perturbation theory and finite Markov chains, \emph{J. Applied Probability} 5 (1968), 401--413.
\end{thebibliography}
\end{document}
